\par\bigskip
\hypertarget{native-379-parabolic-remarks}{}\label{native-379-parabolic-remarks}
\pdfbookmark[2]{Remarks on parabolic motion}{parabolic-bemerkungen}
\begin{center}\textsc{Remarks.}\end{center}
\begingroup\small
\noindent\hspace*{\parindent}In the first volume of »Zeitschrift für Astronomie, herausgegeben von \textsc{Lindenau} und \textsc{Bohnenberger}«, the introduction drafted by \textsc{Lindenau} in November 1815 contains on p. 46 the remark »\textsc{Gauss} has recently reworked nearly the whole theory of comets and drawn up new tables, and we can expect an original work by him on the subject that will likely leave nothing to be desired.« The final words are a turn of phrase frequently used by \textsc{Gauss}, so they were probably taken from a letter from \textsc{Gauss} to \textsc{Lindenau}. This is consistent with the letter to \textsc{Olbers} of 13 Jan. 1815, p. 335, where the same turn of phrase occurs. In a letter to \textsc{Schumacher} of 21 Nov. 1825, \textsc{Gauss} says, »The wish I have always had in my work, to bring it to such a degree of perfection that nothing further could be desired, of course makes it extraordinarily difficult for me.« See also \textsc{Schumacher}'s reply of 2 Dec. 1825. That \textsc{Gauss} worked on determining the orbits of comets especially in December 1814 and January 1815 is also evident from the correspondence with \textsc{Olbers} (pp. 332 and 350), while the paper on this subject printed in Vol. VI, p. 25 ff., had already been presented to the Göttingen Society in September 1813.
\par
\SourceLeaf{380}
\noindent\hspace*{\parindent}Pages 325–370 collect all the notes on this subject found among his papers. Notes [1] and [2] are in notebooks, [3] is on a cover page of a copy of »\textsc{Lambert}, insigniores orbitae cometarum proprietates, Augustae Vindelicorum 1761«, and [4] was compiled from entries in the aforementioned book and a notebook. The footnote on p. 336 comes from the wrapper of a copy of the German translation of the aforementioned paper printed in Vol. VI, p. 25 ff. This copy also contains many other notes on the subject (e.g. a derivation of the formula \(\dfrac{\text{Sector}}{\text{Triangle}}=\dfrac{2+\cos\varphi}{3\cos\varphi}\mathpunct{\text{\raisebox{.45ex}{,}}}\) different from the one given on pp. 329–330), but their content is covered by the printed notes and passages from letters. Besides the table printed on p. 331, there are indications of further auxiliary tables that \textsc{Gauss} intended to draw up to solve the equations that arise.
\par
\noindent\hspace*{\parindent}Curiously, little of the ingenuity employed in the aforementioned notes reappears in note [7], which seems to date from January 1815 and probably contains the »model example« mentioned in the letter to \textsc{Olbers} of 13 January 1815 (p. 337). This note is in a notebook. The notation is readily understood, although it differs in part from that of the paper of September 1813. A computational error also slipped into this note. \textsc{Gauss} had found (cf. p. 341) the value of \(Q\) to be \(272^0\ 40'\ 21''\) instead of \(294^0\ 31'\ 57''\), and later noticed the error himself. He had intended to continue calculating the example with \(P=Q\), because this choice of \(P\) gives the greatest precision (cf. p. 341). Owing to the computational error, however, the rest of the calculation was performed with \(P=272^0\ 40'\ 21''\) (p. 342) and was retained in that form in this edition, since \(P\) can be chosen arbitrarily and the departure from the correct value of \(Q\) is not very substantial.
\par
\noindent\hspace*{\parindent}For the table of parabolic motion printed on pp. 351–367, compare the footnote in the Theoria motus on p. 56 of this volume. As Gauss himself states (p. 353), he calculated this table to 5 decimal places of the arcsecond and intended to print it at intervals of \(\frac1{1000}\) in the argument \(\log(At)\) and to 2 decimal places of the arcsecond. For this reason, the 3rd decimal place has been retained in this edition, so that the interpolation (cf. p. 371) can be performed with the precision \textsc{Gauss} wanted for the interval he intended.
\par
\noindent\hspace*{\parindent}The table is found in a notebook in the form in which it is printed, and in part also in a bundle of loose sheets. In the former it bears the inscription »Completed Feb. 25, 1816«, and in the latter »Completed 1816 Feb. 21.«
\par
\noindent\hspace*{\parindent}The letter printed after the table, which was found among the papers, was evidently addressed to \textsc{Encke}, who helped with the calculations here. It is of particular interest because it shows the whole technique that \textsc{Gauss} used to calculate his tables.
\par
\noindent\hspace*{\parindent}It is striking, moreover, that the article published in 1843 and printed on p. 371 no longer makes any mention of \textsc{Gauss}'s own table.
\par\smallskip
\hfill\textsc{Brendel}.
\par\endgroup\bigskip
\begin{center}\rule{.27\linewidth}{.4pt}\end{center}
